\documentclass[11pt]{article}

\usepackage[margin=2.5cm]{geometry}
\usepackage{amsmath}
\usepackage{booktabs}
\usepackage{siunitx}
\usepackage[hidelinks]{hyperref}

\newcommand{\dg}{^{\circ}}

\title{Stage 2 --- Building the Elbow-Flexion OpenSim Model\\[2pt]
\large \texttt{arm26\_paper\_loaded\_brd\_elbow\_research.osim}}
\author{Vision-Based Optical Simulation}
\date{}

\begin{document}
\maketitle

\noindent
This document describes how the Stage-2 model is \emph{built}: an OpenSim arm26
variant extended for repetitive loaded elbow flexion with a validated
brachioradialis path and a physically represented \SI{2}{\kilogram} hand load,
with a fixed/stabilized shoulder. (Validation, usage and limitations are covered
elsewhere; here the focus is construction only.)

%==============================================================
\section{Purpose and Lineage}

The model turns Stage-1 joint-angle motion into Stage-2 biomechanical structure:
elbow inverse-dynamics torque, muscle moment arms, and the muscle forces/activations
solved by Static Optimization under a real dumbbell load. It is built on stock
arm26 and extended in stages:
\[
\text{stock arm26}
\;\rightarrow\; \texttt{arm26\_paper\_loaded}
\;\rightarrow\; \texttt{arm26\_paper\_loaded\_brd}
\;\rightarrow\; \texttt{..\_brd\_elbow\_research}.
\]

\begin{table}[h]
\centering
\begin{tabular}{ll}
\toprule
\textbf{Build stage} & \textbf{What was added} \\
\midrule
\texttt{arm26\_paper\_loaded} & welded \SI{2}{\kilogram} dumbbell body \\
\texttt{arm26\_paper\_loaded\_brd} & + validated brachioradialis path \\
\texttt{..\_brd\_elbow\_research} & shoulder muscles removed (elbow-focused build) \\
\bottomrule
\end{tabular}
\caption{Incremental construction of the model.}
\end{table}

\paragraph{Provenance.} \textit{Project engineering}, built on the OpenSim arm26
example. The arm26 base and its parameters come from Holzbaur, Murray \& Delp
(2005); OpenSim itself from Delp et al.\ (2007) and Seth et al.\ (2011, 2018).

%==============================================================
\section{Mechanical Structure}

\begin{table}[h]
\centering
\begin{tabular}{ll}
\toprule
\textbf{Property} & \textbf{Value} \\
\midrule
coordinates & 2 \\
bodies & 4 (\texttt{base}, \texttt{r\_humerus}, \texttt{r\_ulna\_radius\_hand}, \texttt{load\_2kg}) \\
muscles & 7 \\
gravity & $(0,\,-9.8066,\,0)\ \si{\meter\per\second\squared}$ \\
\bottomrule
\end{tabular}
\end{table}

The two coordinates are \texttt{r\_shoulder\_elev} (held fixed at about $20\dg$)
and \texttt{r\_elbow\_flex} (driven by the Stage-1 minimum-jerk curl). The
Stage-1 motion file therefore needs only the columns
\texttt{time}, \texttt{r\_shoulder\_elev}, \texttt{r\_elbow\_flex}.

\paragraph{Provenance.} The two-DOF skeleton, joints and the six stock muscles
are \textit{from} the arm26 distribution (Holzbaur, Murray \& Delp 2005; OpenSim,
Delp et al.\ 2007). The fourth body (\texttt{load\_2kg}) is a \textit{project}
addition (Section~\ref{sec:load}).

%==============================================================
\section{Muscle Set}

\begin{table}[h]
\centering
\begin{tabular}{lll}
\toprule
\textbf{Group} & \textbf{Muscles} & \textbf{Role} \\
\midrule
elbow flexors & \texttt{BIClong}, \texttt{BICshort}, \texttt{BRA}, \texttt{BRD\_hand} & curl producers \\
elbow extensors & \texttt{TRIlong}, \texttt{TRIlat}, \texttt{TRImed} & antagonists \\
actuators & \texttt{shoulder\_assist}, \texttt{elbow\_assist} & residual / stabilization \\
\bottomrule
\end{tabular}
\caption{Seven Thelen2003 muscles (4 flexors + 3 extensors) plus two coordinate actuators.}
\end{table}

Six of the seven muscles (the two biceps heads, brachialis, and three triceps
heads) are the stock arm26 set; the seventh, \texttt{BRD\_hand}, is the added
brachioradialis (Section~\ref{sec:brd}).

\paragraph{Provenance.} The six stock muscles and the Thelen2003 force model are
\textit{from} Holzbaur, Murray \& Delp (2005) and Thelen (2003). The
\texttt{shoulder\_assist}/\texttt{elbow\_assist} coordinate actuators are a
\textit{project} addition for residual support (a standard OpenSim practice).

%==============================================================
\section{Brachioradialis: Geometry and Physiology}
\label{sec:brd}

Stock arm26 omits the brachioradialis, the largest elbow-flexor moment arm. It
was added as \texttt{BRD\_hand} and built to behave like human brachioradialis,
not a duplicate brachialis.

\paragraph{Path geometry (the part that sets the moment arm).}
\begin{table}[h]
\centering
\begin{tabular}{lll}
\toprule
\textbf{Point} & \textbf{Frame} & \textbf{Location (m)} \\
\midrule
P1 (origin) & \texttt{r\_humerus} & $(0.0088,\,-0.1806,\,0.0004)$ --- lateral supracondylar ridge \\
P2 (bowstring) & \texttt{r\_ulna\_radius\_hand} & $(0.0234,\,-0.0744,\,0.0440)$ --- anterolateral forearm \\
P3--P6 & \texttt{r\_ulna\_radius\_hand} & route the tendon to the radial styloid \\
\bottomrule
\end{tabular}
\end{table}

Because the origin (P1) is on the humerus and all other points are rigid on the
forearm, only the crossing segment P1$\to$P2 changes length with elbow angle, so
those two points set the elbow moment arm; P3--P6 carry the tendon to the
insertion without changing it.

\paragraph{Physiology (Thelen2003 parameters).}
\begin{table}[h]
\centering
\begin{tabular}{ll}
\toprule
\textbf{Parameter} & \textbf{Value} \\
\midrule
peak isometric force $F_0$ & \SI{261.3}{\newton} \\
optimal fiber length $L_{\text{opt}}$ & \SI{0.173}{\meter} \\
pennation angle & $0\dg$ \\
tendon slack length & \SI{0.165}{\meter} (calibrated) \\
\bottomrule
\end{tabular}
\end{table}

The geometry was tuned so the elbow-flexion moment arm reaches the cadaveric
target ($\SIrange{7.0}{9.0}{\centi\meter}$ peak near $100\dg$--$118\dg$), and the
tendon slack length was calibrated to centre the fiber over the arm26 path
length.

\paragraph{Provenance.} The physiology values are \textit{from} Holzbaur, Murray
\& Delp (2005) / Murray, Buchanan \& Delp (2000) ($F_0$, $L_{\text{opt}}$,
pennation). The moment-arm target is \textit{from} Murray, Buchanan \& Delp
(2002). The reliance on path points/wrap geometry to set moment arms follows
Garner \& Pandy (2000). The specific arm26-frame coordinates and the tendon-slack
calibration are a \textit{project} implementation.

%==============================================================
\section{Removing the Shoulder Muscles}

The four shoulder muscles (\texttt{DELT\_ant}, \texttt{DELT\_post}, \texttt{PECT},
\texttt{LAT}) were removed in the elbow-focused build because:
\begin{itemize}\itemsep2pt
  \item none of them cross the elbow, so they do not affect elbow-flexion loading;
  \item their simplified paths were non-physiological (e.g.\ \texttt{LAT} operated
        at $\sim0.29\times L_{\text{opt}}$), which could contaminate Static Optimization;
  \item stock arm26 has no shoulder muscles either, so removal returns to that
        clean baseline plus the validated elbow extensions.
\end{itemize}
The fixed shoulder is instead held by the \texttt{shoulder\_assist} actuator. The
elbow muscle mechanics are left unchanged.

\paragraph{Provenance.} \textit{Project design decision.} It is motivated by the
stock-arm26 baseline (no shoulder muscles; Holzbaur et al.\ 2005) and by avoiding
a non-physiological optimisation; it is not prescribed by a specific paper. (The
deliberately rejected alternative --- editing one physiology number to mask the
issue --- would not be a valid biomechanical fix.)

%==============================================================
\section{The \SI{2}{\kilogram} Dumbbell Load Body}
\label{sec:load}

The hand load is a real rigid body, not an injected torque.

\begin{table}[h]
\centering
\begin{tabular}{ll}
\toprule
\textbf{Property} & \textbf{Value} \\
\midrule
body & \texttt{load\_2kg} \\
mass & \SI{2.000}{\kilogram} \\
mass centre & $(0,0,0)$ (at the grip) \\
inertia & $(0.0087,\,0.0009,\,0.0087,\,0,0,0)\ \si{\kilogram\meter\squared}$ \\
weld translation & $(0.028,\,-0.285,\,0.082)\ \si{\meter}$ on \texttt{r\_ulna\_radius\_hand} \\
weld rotation & $(\pi/2,\,0,\,0)\ \si{\radian}$ (handle across the palm) \\
\bottomrule
\end{tabular}
\end{table}

The body is welded to the forearm/hand segment; gravity then acts on it and the
elbow load \emph{emerges} from the geometry. The handle axis is oriented across
the palm and nearly parallel to the elbow flexion axis, so it stays steady
through the curl. The inertia tensor is that of a compact dumbbell (a handle plus
two end plates), small about the bar axis and larger across it. The resulting
gravitational elbow torque is consistent with the simple estimate
\[
\tau \approx m\,g\,d \approx \SI{2}{\kilogram}\times\SI{9.81}{\meter\per\second\squared}\times\SI{0.29}{\meter}
\approx \SI{5.7}{\newton\meter}.
\]

\paragraph{Provenance.} The \emph{welded-body} approach to a held load is
\textit{from} the upper-limb literature (Gastaldi et al.\ 2021; Wu et al.\ 2016).
The specific mass, inertia tensor, grip placement and handle orientation are a
\textit{project} implementation.

%==============================================================
\section*{References}
\begin{description}\itemsep2pt
\item[Delp et al.\ (2007).] OpenSim: open-source software to create and analyze
  dynamic simulations of movement. \emph{IEEE TBME}. The simulation platform.
\item[Seth et al.\ (2011, 2018).] OpenSim modeling/simulation framework and
  4.x workflow. \emph{Procedia IUTAM}; \emph{PLOS Comput.\ Biol.} Model
  architecture and API-based construction.
\item[Holzbaur, Murray \& Delp (2005).] A model of the upper extremity\ldots
  \emph{Annals of Biomedical Engineering}. arm26 anatomical lineage and muscle
  parameters.
\item[Thelen (2003).] Adjustment of muscle mechanics model parameters\ldots
  \emph{J.\ Biomech.\ Eng.} The Thelen2003 muscle model used for all muscles.
\item[Garner \& Pandy (2000).] The obstacle-set method for representing muscle
  paths. \emph{CMBBE}. Path-geometry/wrapping basis for moment arms.
\item[Murray, Buchanan \& Delp (2000).] The isometric functional capacity of
  muscles that cross the elbow. \emph{J.\ Biomech.} Elbow-muscle capacity and BRD
  physiology.
\item[Murray, Buchanan \& Delp (2002).] Scaling of peak moment arms of elbow
  muscles\ldots \emph{J.\ Biomech.} Target BRD peak moment arm ($\sim$7--9 cm).
\item[Gastaldi et al.\ (2021); Wu et al.\ (2016).] Loaded upper-limb OpenSim
  modeling with held weights. Basis for the welded-load approach.
\end{description}

\end{document}
